\section{Phase 2: Property-Guided and Context-Aware Composition}

Phase~2 extends compositional generative modelling beyond unconditional
hybridisation and investigates whether \textit{property-guided} sampling can
improve the semantic reliability and clinical plausibility of composed
distributions. This phase directly supports \textbf{Research Goal~2} and tests
\textbf{Hypothesis~H2}, which states that meaningful property-guided
compositional generation requires guidance models that are smooth,
well-calibrated, and robust to distribution shift. In particular, models must
avoid producing out-of-distribution (OOD) pathological structures when steering
composed samples toward continuous clinical attributes such as lesion severity
or spatial involvement.

Where Phase~1 evaluates whether composition alone can produce coherent hybrid
pathology, Phase~2 asks a more refined question: when we condition the composed
distribution on a continuous attribute, does guided sampling remain anatomically
and semantically plausible? This question is crucial for medical imaging, where
severity, extent, and localisation of disease are fundamental descriptors, and
where naive guidance can easily hallucinate unrealistic or clinically dangerous
patterns. Property-guided composition therefore provides a demanding test of
compositional robustness, sensitivity to OOD behaviour, and the structural
conditions needed for reliable continuous attribute control.

\paragraph{Naive Property Guidance.}
We begin by training severity or location predictors on TB, Silicosis, or
combined datasets and applying their gradients to steer samples from the
composed models. This allows us to determine whether gradient-based guidance
interacts coherently with the composed density. Failure modes—such as
exaggerated lesions, anatomically implausible cavity shapes, or misplaced
nodules—indicate that the guidance model is overconfident, insufficiently
smooth, or misaligned with the composed latent space. These observations offer
direct evidence for or against Hypothesis~H2 by revealing whether naive guidance
can preserve semantic meaning or destabilises composition.

\paragraph{Context-Guided and Uncertainty-Aware Predictors.}
To test the second part of Hypothesis~H2, we train context-aware, uncertainty-
calibrated guidance models that incorporate additional latent or image features
before producing severity or localisation predictions. The goal is to ensure
that the guidance model remains reliable when operating on composed samples,
which differ from pure TB or pure Silicosis images. By comparing these models
to naive predictors, we evaluate whether improved calibration and contextual
conditioning mitigate OOD hallucinations and yield more stable control during
sampling. Success indicates that property-guided composition requires guidance
models that explicitly account for uncertainty, distribution shift, and latent
structure. Failure highlights the structural gaps between composition and
attribute control.

\paragraph{Guided Composition Evaluation.}
Across both naive and context-guided approaches, we evaluate the clinical
plausibility of property-steered composite samples by examining whether the
resulting images respect anatomical boundaries, generate realistic lesion
patterns, and maintain coherence with known co-morbidity presentations. The
goal is to determine whether guided operators reflect genuine continuous factors
of variation—such as increasing cavitation severity—or whether guidance drives
the model into regions of latent space that do not correspond to real
pathology. Observing whether composition and guidance reinforce or conflict with
each other provides insight into the structural compatibility of these two
processes.

\paragraph{Contribution of This Phase.}
Phase~2 contributes a unified framework for evaluating how property guidance
interacts with compositional generative modelling. It clarifies when guided
sampling can enhance semantic meaning and when it destabilises composition,
particularly in clinically complex domains. By contrasting naive guidance with
context-aware, uncertainty-calibrated predictors, we identify the structural
conditions under which property-guided composition preserves anatomical
plausibility and yields interpretable continuous variation. This phase therefore
provides a rigorous test of Hypothesis~H2 and establishes the methodological and
theoretical groundwork for integrating continuous attribute control into
compositional generative models.
